\section{模型训练}

\subsection{训练目标与经验风险}
模型训练的目的，是从给定的候选模型中选出一个能够在未知样本上取得良好预测表现的模型。为此，首先需要根据具体任务明确模型表现的评价标准。例如，在分类任务中，可以关注总体分类准确率，也可以更重视某一类别的召回率、精确率或者误分类成本。不同的评价标准体现了不同的任务需求，并决定了哪些预测结果应被认为更好、哪些预测错误应受到更大的惩罚。\par
评价指标主要用于衡量训练完成后的模型是否满足任务需求，但它未必适合直接作为训练时的优化目标。许多评价指标是不连续的、不可微的，或者不能分解为各个样本损失的平均。例如，分类准确率仅取决于最终预测类别是否正确。当模型的预测分数发生小幅变化但预测类别保持不变时，准确率也保持不变，因此它通常无法提供参数应当如何更新的信息。\par
为便于优化，训练时通常使用一个对模型输出变化更加敏感的损失函数作为代理目标。损失函数不必与评价指标具有相同的表达式，但应当与任务目标相容，即损失的减小通常能够促进所关心的预测性能改善。因此，评价指标决定模型最终应达到的目标，损失函数则将这一目标转化为可以进行参数优化的数学问题。\par
设观测数据集$\mathcal{D}=\{(x_i,y_i)\}_{i=1}^{n}$由分布$P$中的$n$个独立同分布样本构成，候选模型组成函数族：
\begin{equation*}
	\mathcal{F}=\{f_\theta:X\to\mathcal{A}:\theta\in\Theta\}
\end{equation*}
其中$\mathcal{A}$为模型输出空间，$\Theta$为参数空间。沿用点估计中的记号，记$L:Y\times\mathcal{A}\to[0,+\infty]$为损失函数，则$L[y,f_\theta(x)]$表示模型$f_\theta$在样本$(x,y)$上产生的预测损失。若分布$P$已知，则模型$f_\theta$在来自同一分布的样本上的平均预测损失可以表示为：
\begin{equation*}
	R(\theta)\coloneq\operatorname{E}(L)
\end{equation*}
称$R(\theta)$为模型$f_\theta$的总体风险。总体风险越小，表示模型在未知样本上产生的平均预测损失越小。因此，从预测角度看，理想的参数应满足：
\begin{equation*}
	\theta^\ast\in\underset{\theta\in\Theta}{\arg\min}R(\theta)
\end{equation*}
然而，分布$P$通常未知，因此总体风险$R(\theta)$无法直接计算。训练时只能利用观测数据，以样本平均损失近似总体风险。
\begin{definition}
	称：
	\begin{equation*}
		\hat{R}(\theta)\coloneq\frac{1}{n}\sum_{i=1}^{n}L[y_i,f_\theta(x_i)]
	\end{equation*}
	为模型$f_\theta$在观测数据集$\mathcal{D}$上的\gls{EmpiricalRisk}。
\end{definition}
以经验风险作为训练目标时，模型训练可以表示为
\begin{equation*}
	\hat{\theta}\in\underset{\theta\in\Theta}{\arg\min}\hat{R}(\theta)
\end{equation*}
这种通过最小化经验风险选择模型参数的方法称为经验风险最小化。由于$\hat{R}(\theta)$依赖于随机抽取的观测数据，因此所得参数$\hat{\theta}$也是数据的函数。
\subsubsection{例子}
下面以分类任务为例，说明如何从模型希望达到的预测效果出发构造损失函数。设类别集合为$Y=\{1,2,\dots,K\}$，分类模型对输入$x$输出一个概率向量：
\begin{equation*}
	q_\theta(x)=(q_{\theta1}(x),q_{\theta22}(),\dots,q_{\theta K}(x))^{\top},\quad
	q_{\theta k}(x)>0,\quad\sum_{k=1}^{K}q_{\theta k}(x)=1
\end{equation*}
其中，$q_{\theta k}(x)$表示模型认为样本$x$属于第$k$类的概率。对于训练样本$(x_i,y_i)$，我们希望模型预测的准确率高，也即赋给真实类别$y_i$的概率$q_{\theta y_i}(x_i)$尽可能接近$1$。\par
由于训练样本相互独立，模型在所有样本上为真实类别给出的概率可以合并为乘积：
\begin{equation*}
	Q_{\mathcal{D}}(\theta)\coloneq\prod_{i=1}^{n}q_{\theta y_i}(x_i)
\end{equation*}
这个乘积越大，说明模型为观测到的真实类别整体赋予的概率越大。最大化$Q_{\mathcal{D}}(\theta)$等价于最小化：
\begin{equation*}
	-\frac{1}{n}\ln Q_{\mathcal{D}}(\theta)=-\frac{1}{n}\sum_{i=1}^{n}\ln q_{\theta y_i}(x_i)
\end{equation*}
由此可以把单个样本的损失定义为：
\begin{equation*}
	L[y_i,q_\theta(x_i)]\coloneq-\ln q_{\theta y_i}(x_i)
\end{equation*}
称该损失为逐样本的\gls{CrossEntropy}损失。当$q_{\theta y_i}(x_i)$接近$1$时，该损失接近$0$；当$q_{\theta y_i}(x_i)$接近$0$时，该损失趋于$+\infty$。因此，模型不仅会因预测类别错误而受到惩罚，还会因确信程度不足而产生较大的损失。相应的经验风险为：
\begin{equation*}
	\hat{R}(\theta)\coloneq\frac{1}{n}\sum_{i=1}^{n}L[y_i,q_\theta(x_i)]=-\frac{1}{n}\sum_{i=1}^{n}\ln q_{\theta y_i}(x_i)
\end{equation*}
%\begin{definition}
%	设 $(X,\mathscr{A},P)$ 为概率空间，$f$为$X$上取值于有限集合
%	$\{x_1,\dots,x_n\}$ 的离散型随机变量，其真实分布为：
%	\begin{equation*}
%		p_i=P(f=x_i),\;i=1,\dots,n
%	\end{equation*}
%	另设$q=(q_1,\dots,q_n)$为定义在同一取值集合上的另一概率分布，满足$q_i>0$。称：
%	\begin{equation*}
%		H(p,q)=-\sum_{i=1}^np_i\ln q_i
%	\end{equation*}
%	为分布$p$相对于分布$q$的\gls{CrossEntropy}。
%\end{definition}
%\begin{property}\label{prop:CrossEntropy}
%	设$p=(p_1,\dots,p_n)$与$q=(q_1,\dots,q_n)$（$p_iq_i\ne0,\;i=1,2,\dots,n$）为定义在同一有限集合上的概率分布，则$H(p,q)\geqslant H(p,p)$，上式取等当且仅当$q_1=p_1,q_2=p_2,\dots,q_n=p_n$ 。
%\end{property}
%\begin{proof}
%	因为：
%	\begin{equation*}
%		H(p,q)=-\sum_{i=1}^np_i\ln q_i=-\sum_{i=1}^np_i\ln\frac{q_i}{p_i}-\sum_{i=1}^np_i\ln p_i
%	\end{equation*}
%	并且由\cref{prop:ConvexFunction}(3)可知$-\ln(x)$在$(0,+\infty)$上为凸函数，由\cref{ineq:Jensen}可知：
%	\begin{equation*}
%		\sum_{i=1}^np_i\ln\frac{q_i}{p_i}\leqslant\ln\left[\sum_{i=1}^{n}p_i\frac{q_i}{p_i}\right]=0
%	\end{equation*}
%	上式取等当且仅当$\dfrac{q_i}{p_i}$为常数，即$q_i=p_i$。所以：
%	\begin{equation*}
%		H(p,q)\geqslant-\sum_{i=1}^np_i\ln p_i=H(p,p)
%	\end{equation*}
%	且当且仅当$q=p$时取等。
%\end{proof}
%交叉熵需要区分真实概率$p$和预测概率$q$。固定真实概率$p=0.25$后，预测概率$q$与$p$一致时交叉熵最小；预测越偏离真实概率，交叉熵越大。
%\begin{figure}[H]
%	\centering
%	\begin{tikzpicture}
%		\begin{axis}[
%			width=9.2cm, height=5.6cm,
%			xlabel={$q$}, ylabel={$H(p,q)$},
%			xmin=0, xmax=1, ymin=0, ymax=4.5,
%			xtick={0,0.25,0.5,0.75,1}, ytick={0,1,2,3,4},
%			axis lines=left, grid=major,
%			grid style={draw=gray!18},
%			samples=200,
%			]
%			\addplot[very thick, orange!85!black, domain=0.02:0.98]
%			{-(0.25*ln(x)+0.75*ln(1-x))};
%			\addplot[only marks, mark=*, orange!85!black] coordinates {(0.25,0.56234)};
%			\draw[dashed, gray] (axis cs:0.25,0) -- (axis cs:0.25,0.56234);
%		\end{axis}
%	\end{tikzpicture}
%	\caption{真实概率固定为$p=0.25$时的二分类交叉熵。横轴$q$为模型预测第一类的概率。}
%\end{figure}
%\begin{note}
%	由\cref{prop:CrossEntropy}可知，交叉熵在所有概率分布$q$中以$q=p$为唯一最小值。实际中真实分布$p$通常未知，故以样本的经验分布$\hat{p}$代替$p$，并最小化$H(\hat{p},q)$。在$H(\hat{p},q)$能够充分逼近$H(p,q)$的条件下，该过程使所估计的分布$q$趋近于真实分布$p$。
%\end{note}
%\begin{note}
%	在分类问题中，我们可以用逐样本的交叉熵作为损失函数。对于第$i$个样本，若其真实类别为第$k$类，则其真实分布$p_i$是一个one-hot向量，即第$k$个分量为$1$，其余分量为$0$。例如，在二分类问题中有：
%	\begin{equation*}
%		p_i=(1,0)^{\top}\qquad\text{或}\qquad p_i=(0,1)^{\top}
%	\end{equation*}
%	若模型对该样本输出概率分布$q_i$，则其交叉熵为：
%	\begin{equation*}
%		H(p_i,q_i)=-\sum_{j=1}^Kp_{ij}\ln q_{ij}=-\ln q_{iy_i}
%	\end{equation*}
%	整个数据集上的交叉熵损失通常取各样本损失的平均：
%	\begin{equation*}
%		L_{\mathrm{CE}}=\frac{1}{n}\sum_{i=1}^nH(p_i,q_i)
%	\end{equation*}
%	由此可见，分类问题中的交叉熵分别比较每个样本的真实分布$p_i$与预测分布$q_i$。由于$p_i$是one-hot向量，$H(p_i,q_i)$实际上只取决于模型赋给该样本真实类别的概率$q_{iy_i}$：该概率越接近$1$，损失越小；该概率越接近$0$，损失越大。
%\end{note}

\subsection{过拟合}
经验风险$\hat{R}(\theta)$只度量模型在观测数据集上的平均损失，总体风险$R(\theta)$则度量模型对同一分布中新样本的平均损失，二者通常并不相等。\par
有时持续降低$\hat{R}(\theta)$未必会同时降低$R(\theta)$：模型可能不是在提取能够推广到新样本的稳定规律，而是利用观测数据中的偶然噪声、异常点甚至样本的细微特征，对有限训练集作近似记忆。此时训练损失较低，而新样本上的损失较高，这种现象称为\gls{Overfitting}。\par
\begin{figure}[htbp]
	\centering
	\makebox[\textwidth][c]{%
		\begin{tikzpicture}
			\begin{groupplot}[
				group style={
					group size=3 by 1,
					horizontal sep=0.025\textwidth
				},
				width=0.30\textwidth,
				height=0.27\textwidth,
				xmin=-1.25,
				xmax=1.25,
				ymin=-2.2,
				ymax=2.2,
				domain=-1.2:1.2,
				samples=300,
				axis lines=middle,
				xlabel={$x$},
				ylabel={$y$},
				xtick={-1,0,1},
				ytick={-2,-1,0,1,2},
				tick label style={font=\small},
				label style={font=\small},
				title style={
					font=\small,
					yshift=-0.5ex
				},
				clip=false
				]
				
				% 第一种复杂拟合
				\nextgroupplot[
				title={Cubic fit}
				]
				
				\addplot[
				dashed,
				thick,
				gray
				] {x};
				
				\addplot[
				thick
				] {
					x
					+1.8*(x+1)*x*(x-1)
				};
				
				\addplot[
				only marks,
				mark=*,
				mark size=2.2pt
				] coordinates {
					(-1,-1)
					(0,0)
					(1,1)
				};
				
				% 第二种复杂拟合
				\nextgroupplot[
				title={Quintic fit},
				legend style={
					at={(0.5,-0.34)},
					anchor=north,
					legend columns=-1,
					draw=none,
					font=\small,
					/tikz/every even column/.append style={column sep=0.8em}
				}
				]
				
				\addplot[
				dashed,
				thick,
				gray
				] {x};
				\addlegendentry{true relation $y=x$}
				
				\addplot[
				thick
				] {
					x
					-3.2*(x+1)*x*(x-1)*(x^2-0.35)
				};
				\addlegendentry{fitted curve}
				
				\addplot[
				only marks,
				mark=*,
				mark size=2.2pt
				] coordinates {
					(-1,-1)
					(0,0)
					(1,1)
				};
				\addlegendentry{training point}
				
				% 第三种复杂拟合
				\nextgroupplot[
				title={Oscillatory fit}
				]
				
				\addplot[
				dashed,
				thick,
				gray
				] {x};
				
				\addplot[
				thick
				] {
					x
					+1.1*(x+1)*x*(x-1)
					*sin(deg(7*x))
				};
				
				\addplot[
				only marks,
				mark=*,
				mark size=2.2pt
				] coordinates {
					(-1,-1)
					(0,0)
					(1,1)
				};
				
			\end{groupplot}
			
		\end{tikzpicture}%
	}
	\caption{取自$y=x$的三个样本可由不同的模型进行解释。}
	\label{fig:overfitting-three-points}
\end{figure}
从\cref{fig:overfitting-three-points}可以看出，三条复杂程度不同的曲线都完美地拟合了由三个观测点构成的观测数据集，因此它们在该观测集上取得了相同且最小的经验风险$\hat{R}(\theta)$。然而，在训练样本之外，若新的观测仍来自真实关系$y=x$，这些模型会产生较大的预测误差，即它们存在过拟合的现象。由此可见，仅凭观测数据上的经验风险无法判断模型的真实预测能力；即使多个模型具有相同的训练误差，它们的总体风险$R(\theta)$也可能存在显著差异。\par
综上，我们需要关注差值：
\begin{equation*}
	R(\theta)-\hat{R}(\theta)
\end{equation*}
它反映了模型从观测数据推广到新数据时的性能差异。
\subsubsection{训练集、验证集与测试集}
计算上述差值需要计算$R(\theta)$，但是由于分布$P$未知，$R(\theta)$通常无法直接计算，只能通过未参与训练的数据加以估计，所以我们通常将原始数据随机划分为三个互不相交的部分：
\begin{equation*}
	\mathcal{D}=\mathcal{D}_{\mathrm{train}}\cup\mathcal{D}_{\mathrm{val}}\cup\mathcal{D}_{\mathrm{test}}
\end{equation*}
分别称为\gls{TrainingSet}、\gls{ValidationSet}和\gls{TestSet}。三者应当尽可能来自相同的目标分布；对于类别不平衡的分类问题，通常采用分层抽样，使各部分的类别比例大致一致。\par
\begin{method}
	数据集划分后的模型训练与选择通常按下列顺序进行：
	\begin{enumerate}
		\item 在训练集$\mathcal{D}_{\mathrm{train}}$上估计模型的可学习参数$\theta$；
		\item 在验证集$\mathcal{D}_{\mathrm{val}}$上比较不同模型的性能，并据此选择模型；
		\item 确定全部建模方案后，只在测试集$\mathcal{D}_{\mathrm{test}}$上进行一次最终评价，报告模型对未见数据的性能估计。
	\end{enumerate}
\end{method}
验证集虽然不参与参数的直接拟合，却参与模型选择；若反复根据测试集结果修改模型，测试集实际上也参与了训练流程，其评价结果便会过于乐观。因此，测试集必须与整个训练和选择过程隔离。数据预处理也应遵循相同原则：缺失值填补、标准化、特征选择等步骤只能在训练集上估计所需参数，再将已经确定的变换应用到验证集和测试集，不能提前利用后两者的信息。\par
当样本量较小时，如果再从训练数据中单独划分一个固定的验证集会进一步减少模型实际用于训练的样本。此时可以在训练数据内部使用$K$折\gls{CV}：将训练数据近似均匀地分为$K$份，依次将其中一份作为验证集，其余$K-1$份作为训练集，共进行$K$轮训练与验证，并对$K$次验证结果取平均，以评价候选模型。此时仍应保留独立测试集；交叉验证用于模型选择，测试集只用于最终评价。
\subsubsection{正则化与模型空间约束}
对$f=f_\theta\in\mathcal{F}$，记$\hat{R}_{\mathrm{train}}(f)\coloneq\hat{R}_{\mathrm{train}}(\theta)$及$R(f)\coloneq R(\theta)$。\par
若$\mathcal{F}_1\subseteq\mathcal{F}_2$，则：
\begin{equation*}
	\inf_{f\in\mathcal{F}_2}\hat{R}_{\mathrm{train}}(f)\leqslant\inf_{f\in\mathcal{F}_1}\hat{R}_{\mathrm{train}}(f)
\end{equation*}
也就是说，扩大模型空间不会使最优训练风险变大，因为新增的模型为拟合训练数据提供了更多选择；但总体风险并不具有相同的单调性。模型空间越大，训练过程越容易从大量候选模型中找到一个恰好迎合当前样本偶然性的模型，因而经验风险虽然可以很低，经验风险与总体风险之间的偏差却更难控制。极端情况下，一个足够复杂的模型可以记住每个训练样本，使训练风险接近$0$；但训练样本之外的函数取值仍可能几乎不受约束，这正是过拟合产生的原因。
\begin{definition}
	设$\Omega:\mathcal{F}\to[0,+\infty]$为模型复杂度函数，$\Omega(f)$越大表示模型$f$越复杂。对$c\geqslant0$，定义受限模型空间：
	\begin{equation*}
		\mathcal{F}_c\coloneq\{f\in\mathcal{F}:\Omega(f)\leqslant c\}
	\end{equation*}
	在$\mathcal{F}_c$中最小化经验风险：
	\begin{equation*}
		\hat{f}_c\in\underset{f\in\mathcal{F}_c}{\arg\min}\hat{R}_{\mathrm{train}}(f)
	\end{equation*}
	称为约束形式的\gls{Regularization}。其中，$c$称为复杂度上界。
\end{definition}
实际计算中更常使用正则化的惩罚形式：
\begin{equation*}
	\hat{f}_\lambda\in\underset{f\in\mathcal{F}}{\arg\min}\left\{\hat{R}_{\mathrm{train}}(f)+\lambda\Omega(f)\right\},\;\lambda\geqslant0
\end{equation*}
其中，$\hat{R}_{\mathrm{train}}(f)+\lambda\Omega(f)$称为正则化目标函数，$\lambda$称为\gls{RegularizationStrength}。当$\lambda=0$时，该问题退化为原来的经验风险最小化；$\lambda$越大，模型复杂度在目标函数中付出的代价越大，训练过程便越倾向于选择简单模型。
\begin{note}
	限制模型空间的本质是在拟合训练数据与控制模型复杂度避免过拟合之间取得平衡：模型空间过小会导致模型无法很好的拟合训练数据，也即\gls{Underfitting}，模型空间过大会导致过拟合。
\end{note}
